\section{Numerical evidence}\label{sec:numerico}

Nobel, Agrawal and Boyd~\cite{NAB} give $\alpha\in[1.944000752,\,1.944001082]$.
For the two tables we use $\alpha_0=1.944001$ and $\Phi_0=1-\alpha_0$ as reference values,
because $\alpha=-s_1$ (Section~\ref{sec:preliminares}). The extrapolations have no certified error bound,
and the form of the fits, in powers of $1/n$, is a working hypothesis: Theorem~\ref{thm:intro-constante}
gives the limit, but not the rate of convergence.

\paragraph{The predictor.}
Set $C'_{{\rm pred},0}(n)=\log Q_{\rm pred}(n)-\log n!-n\Phi_0$. A numerical solution for the canonical
maximizer
finds it to be interior at all ten tabulated sizes, in floating-point arithmetic.
The pseudodeterminant of $S_n$ completes the calculation in Table~\ref{tab:predictor}.
With $\alpha_0$ and $40\le n\le240$, the fit $C+a/n+b/n^2$ gives $C=0.94015$ and the fit that adds $c/n^3$
gives $C=0.94013$.
Over the full interval for $\alpha$, the fits starting at $40$ or $60$ give rounded endpoints of $0.9400$
and $0.9402$.
When $20$ and $27$ are included, the rounded endpoints are $0.9397$ and $0.9404$. This gives the estimate
$C_{\rm pred}\approx0.9401$.

\begin{table}[ht]
\centering
\begin{tabular}{lccccc}
\toprule
$n$ & 20 & 27 & 40 & 60 & 80\\
$C'_{{\rm pred},0}(n)$ & 0.9112 & 0.9165 & 0.9239 & 0.9290 & 0.9316\\
\midrule
$n$ & 100 & 130 & 160 & 200 & 240\\
$C'_{{\rm pred},0}(n)$ & 0.9333 & 0.9348 & 0.9358 & 0.9367 & 0.9372\\
\bottomrule
\end{tabular}
\caption{Constant term of the predictor for finite boards.}\label{tab:predictor}
\end{table}

\paragraph{Direct evaluation of the formula.}
Separately, we computed the continuous profile $\nu_1$ by solving its equations at mass one numerically. We
used it to compute $I_2$ and $\det_2(I+J)$ by a Nyström discretization~\cite{Bornemann} of the operator in
Section~\ref{sec:D03}, on successive grids with extrapolation in the grid spacing. This gives
$I_2\approx1.02638534$ and $\det_2(I+J)\approx0.005162694$, and the formula in
Theorem~\ref{thm:intro-constante} gives $C_{\rm pred}\approx0.940144$, within the bands above. This
calculation uses neither $\alpha$ nor Table~\ref{tab:predictor}, and it too has no certified error bound.
For the pairs of grids of sizes $192$ and $384$ (even), and $193$ and $385$ (odd), the extrapolations gave
$0.94014395698$ and $0.94014395711$, respectively.

\paragraph{The exact count.}
Exact values of $Q(n)$ up to $n=27$~\cite{OEIS,PE} have been known since 2016; the value for $n=27$ was
confirmed in 2025 with another program~\cite{YaoLi}. For $16\le n\le27$, the values of $C'_0(n)=\log
Q(n)-\log n!-n\Phi_0$ are given in Table~\ref{tab:exacto}.

\begin{table}[ht]
\centering
\begin{tabular}{lcccccc}
\toprule
$n$ & 16 & 17 & 18 & 19 & 20 & 21\\
$C'_0(n)$ & 0.9404 & 0.9209 & 0.9135 & 0.9224 & 0.9320 & 0.9187\\
\midrule
$n$ & 22 & 23 & 24 & 25 & 26 & 27\\
$C'_0(n)$ & 0.9178 & 0.9241 & 0.9295 & 0.9272 & 0.9265 & 0.9285\\
\bottomrule
\end{tabular}
\caption{Constant term of the exact count.}\label{tab:exacto}
\end{table}

In November 2026 a preprint reported a value for $Q(28)$~\cite{Cherepanov},
obtained using a new exact counting algorithm. Shortly afterwards the OEIS~\cite{OEIS} also added
$Q(29)$, computed by P.~A.~Webb with an implementation of the color decomposition that, according to
that entry, was used for $Q(28)$ and $Q(29)$. To our knowledge, neither value has been
verified with a different algorithm, so we do not include them in Table~\ref{tab:exacto} or
Figure~\ref{fig:comparacion}. For information, the first gives $C'_0(28)\approx0.9305$, within
the range of the tabulated values.

Figure~\ref{fig:comparacion} shows the two tables on the same scale.
As $n$ increases, the points approach the vertical axis because the horizontal coordinate is $1/n$.

\begin{figure}[htbp]
\centering
\begin{tikzpicture}[x=158cm,y=103cm]
 \begin{scope}[shift={(0,-.905)}]
 \foreach \v/\letra in {.910/{0.91},.920/{0.92},.930/{0.93},.940/{0.94}}{
   \draw[black!18] (0,\v)--(.065,\v);
   \node[anchor=east,font=\footnotesize] at (-.0007,\v) {\letra};}
 \draw[btinta] (0,.905)--(.066,.905) (0,.905)--(0,.945);
 \foreach \v/\letra in {0/{0},.02/{0.02},.04/{0.04},.06/{0.06}}{
   \draw[btinta] (\v,.905)--(\v,.9045);
   \node[anchor=north,font=\footnotesize] at (\v,.904) {\letra};}
 \draw[btinta,dashed] (0,.940144)--(.065,.940144);
 \draw[btinta,line width=.75pt] plot coordinates
{(0.050000000,0.911188472) (0.037037037,0.916522978) (0.025000000,0.923873790) (0.016666667,0.928966033)
(0.012500000,0.931643280) (0.010000000,0.933288384) (0.007692308,0.934831567) (0.006250000,0.935807595)
(0.005000000,0.936660345) (0.004166667,0.937232189)};
 \foreach \x/\y in {0.050000000/0.911188472,0.037037037/0.916522978,0.025000000/0.923873790,%
0.016666667/0.928966033,0.012500000/0.931643280,0.010000000/0.933288384,0.007692308/0.934831567,%
0.006250000/0.935807595,0.005000000/0.936660345,0.004166667/0.937232189}{\fill[btinta] (\x,%
\y) circle[radius=1.45pt];}
 \foreach \x/\y in {0.062500000/0.940434608,0.058823529/0.920874442,0.055555556/0.913509083,%
0.052631579/0.922429638,0.050000000/0.931979175,0.047619048/0.918660405,0.045454545/0.917777869,%
0.043478261/0.924122583,0.041666667/0.929528249,0.040000000/0.927235036,0.038461538/0.926480702,%
0.037037037/0.928458576}{\draw[btinta,fill=white,line width=.65pt]
    (\x,\y) ++(-1.55pt,-1.55pt) rectangle ++(3.1pt,3.1pt);}
 \node[anchor=north,font=\small] at (.033,.9012) {$1/n$};
 \end{scope}
\end{tikzpicture}
\caption{Centered terms with the same reference $\Phi_0$. The circles, joined by line segments,
represent the predictor in Table~\ref{tab:predictor}; the unconnected squares represent the exact count
in Table~\ref{tab:exacto}. The dashed horizontal line marks the value of the predictor constant
from the Nyström discretization. The sizes shown are those in the two tables.}
\label{fig:comparacion}
\end{figure}


At $n=20$ and $n=27$ we have $\log Q(n)-\log Q_{\rm pred}(n)=0.0208$ and $0.0119$. The two numbers
are compatible with a difference of order $1/n$, but they do not allow us to estimate its limit.

The programs and data used to obtain the numerical values in this section form the numerical verification
package
accompanying the paper~\cite{Codigo}.
